\section{AlphaCode：在程序空间里做大规模搜索}

\subsection{为什么竞赛编程比代码补全难}

HumanEval 常给出函数签名、局部说明与较短实现；Codeforces 问题则要求理解长题意、选择算法、满足复杂度约束并写出完整程序。错误可能来自题意理解、算法、边界条件、复杂度、语法或输入输出协议中的任一环节。

\videofigure{figures/alphacode-pipeline.jpg}{AlphaCode 将代码预训练、竞赛题微调、大规模采样、过滤聚类和最终提交串成完整流水线。}{00:03:48--00:05:42}

训练数据主要来自 GitHub 代码与 CodeContests。原始 AlphaCode 仍需要自建模型：预训练学习通用代码分布，竞赛数据微调让模型适应“题面--解法”结构，之后才进入测试时搜索。

\begin{warningbox}{验证集 loss 不是 solve rate 的可靠代理}
token-level loss 衡量局部预测是否接近训练分布，但竞赛题按整段程序是否通过隐藏测试计分。一个 token 的边界错误就可能让整题归零；因此训练 loss 的小幅改善未必转化为题目通过率。
\end{warningbox}

\subsection{一题采样一百万个程序}

\videofigure{figures/large-scale-sampling.jpg}{AlphaCode 对每道题生成约一百万个 Python/C++ 候选，并用高温度与提示扰动维持算法多样性。}{00:06:08--00:07:28}

设基础模型对程序 $y$ 的条件分布为 $p_\theta(y\mid x)$。若单次采样命中正确程序的概率为 $p$，独立采样 $n$ 次至少命中一次的概率为：

$$
P(\text{hit})=1-(1-p)^n.
$$

\begin{itemize}
    \item $x$：题面、样例和可选标签；
    \item $y$：完整候选程序；
    \item $p$：单次采样落入正确解集合的概率；
    \item $n$：采样预算。
\end{itemize}

当 $p$ 极小时，大规模采样能显著提高覆盖率；但收益最终受选择器限制，因为平台只允许提交很少的程序。

\begin{importantbox}{高温度的目标是覆盖算法族，而非随机制造语法差异}
理想多样性来自动态规划、贪心、图算法、数据结构等不同解题思路。如果候选只是变量名不同或同一错误模板的改写，样本数再大也不增加有效覆盖。
\end{importantbox}

\subsection{先过样例，再按语义聚类}

\videofigure{figures/filter-cluster.jpg}{系统先过滤无法通过公开样例的程序，再用生成测试输入上的行为把语义近似的程序聚成一类。}{00:07:28--00:09:04}

AlphaCode 训练额外的 test-input generator。对同一题生成输入集合 $T=\{t_j\}$，把程序 $y$ 映射为行为签名：

$$
\phi_T(y)=\bigl(y(t_1),y(t_2),\ldots,y(t_m)\bigr).
$$

若两个程序的行为签名相近，它们可能实现同一种算法或共享同一错误；聚类后每簇只保留少数代表，可以在固定提交预算下覆盖更多不同语义。

\begin{knowledgebox}{聚类是一种预算分配器}
平台提交次数是稀缺资源。聚类不是为了让程序看起来整齐，而是避免把十次提交都浪费在同一种失败模式上。
\end{knowledgebox}

\subsection{本章小结}

AlphaCode 的关键是把极大的程序空间压缩成少量多样候选：模型负责 proposal，公开样例做廉价过滤，生成测试做行为聚类，隐藏测试完成最终验证。

